% GENERATED FILE. Edit canonical lessons, then run npm run generate:books.
% canonical-source-hash: fnv1a64:7a0d7729cae4cbbe
% canonical-lessons: LA-C510-ineptus, LA-C510-insanus, LA-C510-insania, LA-C510-demens, LA-C510-deliro

\chapter{Folly and madness}
\label{ch:la-folly-and-madness}
\hlchaptermodality{510}

\begin{chapteropening}
\textbf{By the end of this chapter:} \emph{I can describe foolish and crazy behaviour.}

Complete the last lesson of chapter 510: I can describe foolish and crazy behaviour.
\end{chapteropening}

\section[ineptus, inepta, ineptum]{ineptus, inepta, ineptum — silly, foolish; tactless}

\label{lesson:LA-C510-ineptus}



\begin{quote}
Before the new one: say the Latin for steadfast, then the Latin for good sense.
\end{quote}

\subsection*{You'll want to know: ineptus, inepta, ineptum}
\textbf{ineptus, inepta, ineptum} — \textquotedblleft{}silly, foolish; tactless\textquotedblright{}.

\begin{quote}
Fābula inepta est, sed puerī rīdent.
\end{quote}

\begin{quote}
\emph{The story is silly, but the children laugh.}
\end{quote}

\emph{in-} 'not' + \emph{aptus} 'fitting': someone who says the wrong thing at the wrong time. English \emph{inept}.

\begin{cousinweb}
\textbf{ineptus} is built on \emph{in-} \textquotedblleft{}not\textquotedblright{} and \emph{aptus} \textquotedblleft{}fitting\textquotedblright{} (Lewis \& Short). English \emph{inept} comes from it.
\end{cousinweb}

\begin{grammarlens}[title={What you've built}]
One more word for telling, explaining and describing at B1.
\end{grammarlens}

\subsection*{Your turn}
\begin{itemize}
\raggedright
  \item \emph{Say it:} \emph{ineptus, inepta, ineptum}
  \item \emph{Say it:} \emph{ineptus, inepta, ineptum}, once more
  \item \emph{Say it:} a sentence of your own with \emph{ineptus, inepta, ineptum}
  \item \emph{Recall:} say the Latin for steadfast, then the Latin for good sense, then say \emph{ineptus, inepta, ineptum} again
\end{itemize}

\begin{tcolorbox}[breakable,colback=teal!4,colframe=teal!35!black,title={Before you move on}]
What does ineptus, inepta, ineptum mean? (\textquotedblleft{}Silly, foolish; tactless\textquotedblright{}.) Say it once more.
\end{tcolorbox}
\section[īnsānus, īnsāna, īnsānum]{īnsānus, īnsāna, īnsānum — mad, insane}

\label{lesson:LA-C510-insanus}



\begin{quote}
Before the new one: say the Latin for good sense, then the Latin for silly.
\end{quote}

\subsection*{You'll want to know: īnsānus, īnsāna, īnsānum}
\textbf{īnsānus, īnsāna, īnsānum} — \textquotedblleft{}mad, insane\textquotedblright{}.

\begin{quote}
Vir īnsānus in viā clāmat.
\end{quote}

\begin{quote}
\emph{A madman is shouting in the street.}
\end{quote}

\emph{in-} 'not' + \emph{sānus} 'healthy'. Also of wild, crazy plans. English \emph{insane}.

\begin{cousinweb}
\textbf{īnsānus} is built on \emph{in-} \textquotedblleft{}not\textquotedblright{} and \emph{sānus} \textquotedblleft{}sound, healthy\textquotedblright{} (Lewis \& Short; EtymDB). English \emph{insane} comes from it.
\end{cousinweb}

\begin{grammarlens}[title={What you've built}]
One more word for telling, explaining and describing at B1.
\end{grammarlens}

\subsection*{Your turn}
\begin{itemize}
\raggedright
  \item \emph{Say it:} \emph{īnsānus, īnsāna, īnsānum}
  \item \emph{Say it:} \emph{īnsānus, īnsāna, īnsānum}, once more
  \item \emph{Say it:} a sentence of your own with \emph{īnsānus, īnsāna, īnsānum}
  \item \emph{Recall:} say the Latin for good sense, then the Latin for silly, then say \emph{īnsānus, īnsāna, īnsānum} again
\end{itemize}

\begin{tcolorbox}[breakable,colback=teal!4,colframe=teal!35!black,title={Before you move on}]
What does īnsānus, īnsāna, īnsānum mean? (\textquotedblleft{}Mad, insane\textquotedblright{}.) Say it once more.
\end{tcolorbox}
\section[īnsānia, īnsāniae]{īnsānia, īnsāniae — madness, folly}

\label{lesson:LA-C510-insania}



\begin{quote}
Before the new one: say the Latin for silly, then the Latin for mad.
\end{quote}

\subsection*{You'll want to know: īnsānia, īnsāniae}
\textbf{īnsānia, īnsāniae} — \textquotedblleft{}madness, folly\textquotedblright{}.

\begin{quote}
Haec īnsānia est, nōn cōnsilium.
\end{quote}

\begin{quote}
\emph{This is madness, not a plan.}
\end{quote}

Feminine, first declension, from \emph{īnsānus}.

\begin{cousinweb}
\textbf{īnsānia} is built on \emph{īnsānus} \textquotedblleft{}mad\textquotedblright{} (Lewis \& Short). English \emph{insanity} comes from the related \emph{īnsānitās}.
\end{cousinweb}

\begin{grammarlens}[title={What you've built}]
One more word for telling, explaining and describing at B1.
\end{grammarlens}

\subsection*{Your turn}
\begin{itemize}
\raggedright
  \item \emph{Say it:} \emph{īnsānia, īnsāniae}
  \item \emph{Say it:} \emph{īnsānia, īnsāniae}, once more
  \item \emph{Say it:} a sentence of your own with \emph{īnsānia, īnsāniae}
  \item \emph{Recall:} say the Latin for silly, then the Latin for mad, then say \emph{īnsānia, īnsāniae} again
\end{itemize}

\begin{tcolorbox}[breakable,colback=teal!4,colframe=teal!35!black,title={Before you move on}]
What does īnsānia, īnsāniae mean? (\textquotedblleft{}Madness, folly\textquotedblright{}.) Say it once more.
\end{tcolorbox}
\section[dēmēns, dēmentis]{dēmēns, dēmentis — out of one's mind, crazy}

\label{lesson:LA-C510-demens}



\begin{quote}
Before the new one: say the Latin for mad, then the Latin for madness.
\end{quote}

\subsection*{You'll want to know: dēmēns, dēmentis}
\textbf{dēmēns, dēmentis} — \textquotedblleft{}out of one's mind, crazy\textquotedblright{}.

\begin{quote}
Dēmēns es, sī in tempestāte nāvigās.
\end{quote}

\begin{quote}
\emph{You are out of your mind if you sail in a storm.}
\end{quote}

One ending. \emph{dē-} 'away from' + \emph{mēns} 'mind'. English \emph{demented}.

\begin{cousinweb}
\textbf{dēmēns} is built on \emph{dē-} \textquotedblleft{}away from\textquotedblright{} and \emph{mēns} \textquotedblleft{}mind\textquotedblright{} (Lewis \& Short; EtymDB). English \emph{demented} comes from the related Late Latin verb \emph{dēmentō}.
\end{cousinweb}

\begin{grammarlens}[title={What you've built}]
One more word for telling, explaining and describing at B1.
\end{grammarlens}

\subsection*{Your turn}
\begin{itemize}
\raggedright
  \item \emph{Say it:} \emph{dēmēns, dēmentis}
  \item \emph{Say it:} \emph{dēmēns, dēmentis}, once more
  \item \emph{Say it:} a sentence of your own with \emph{dēmēns, dēmentis}
  \item \emph{Recall:} say the Latin for mad, then the Latin for madness, then say \emph{dēmēns, dēmentis} again
\end{itemize}

\begin{tcolorbox}[breakable,colback=teal!4,colframe=teal!35!black,title={Before you move on}]
What does dēmēns, dēmentis mean? (\textquotedblleft{}Out of one's mind, crazy\textquotedblright{}.) Say it once more.
\end{tcolorbox}
\section[dēlīrō, dēlīrāre]{dēlīrō, dēlīrāre — to be delirious, talk nonsense}

\label{lesson:LA-C510-deliro}



\begin{quote}
Before the new one: say the Latin for madness, then the Latin for out of one's mind.
\end{quote}

\subsection*{You'll want to know: dēlīrō, dēlīrāre}
\textbf{dēlīrō, dēlīrāre} — \textquotedblleft{}to be delirious, talk nonsense\textquotedblright{}.

\begin{quote}
Aeger febrī dēlīrat.
\end{quote}

\begin{quote}
\emph{The patient is delirious with fever.}
\end{quote}

\emph{dēlīrō, dēlīrāre, dēlīrāvī}: \emph{dē} 'out of' + \emph{līra} 'furrow' — to plough crooked, then to wander in mind. English \emph{delirium}.

\begin{cousinweb}
\textbf{dēlīrō} is built on \emph{dē-} \textquotedblleft{}away from\textquotedblright{} and \emph{līra} \textquotedblleft{}a furrow\textquotedblright{} (Lewis \& Short; EtymDB): first to plough out of the line, then to wander in mind. English \emph{delirious} and \emph{delirium} come from it.
\end{cousinweb}

\begin{grammarlens}[title={What you've built}]
One more word for telling, explaining and describing at B1.
\end{grammarlens}

\subsection*{Your turn}
\begin{itemize}
\raggedright
  \item \emph{Say it:} \emph{dēlīrō, dēlīrāre}
  \item \emph{Say it:} \emph{dēlīrō, dēlīrāre}, once more
  \item \emph{Say it:} a sentence of your own with \emph{dēlīrō, dēlīrāre}
  \item \emph{Recall:} say the Latin for madness, then the Latin for out of one's mind, then say \emph{dēlīrō, dēlīrāre} again
\end{itemize}

\begin{tcolorbox}[breakable,colback=teal!4,colframe=teal!35!black,title={Before you move on}]
What does dēlīrō, dēlīrāre mean? (\textquotedblleft{}To be delirious, talk nonsense\textquotedblright{}.) Say it once more.
\end{tcolorbox}
